\documentclass[12pt]{article}
\newcommand{\DocShort}{Ajuste sobre ln(P) con funciones de Excel}
\input{preamble_population.tex}
\begin{document}
\doccover{Ajuste lineal sobre $\ln(P)$ con funciones de Excel}{2026-05}
\handoutintro{Transformar la población con logaritmo natural, ajustar una recta con LINEST y luego deshacer la transformación.}{Cuando se quiere mostrar explícitamente el procedimiento exponencial a partir de una linealización.}

\section{Datos ejemplo}
Aquí se usa $x$ como tiempo medido en años desde 1970.

\begin{center}
\begin{tabular}{lcccccc}
\toprule
Año & 1970 & 1980 & 1990 & 2000 & 2010 & 2020 \\
\midrule
$x$ (años) & 0 & 10 & 20 & 30 & 40 & 50 \\
Población & 26,000 & 29,800 & 34,200 & 39,200 & 45,000 & 51,600 \\
$\ln(P)$ & 10.1659 & 10.3023 & 10.4401 & 10.5767 & 10.7144 & 10.8517 \\
\bottomrule
\end{tabular}
\end{center}

\section{Ecuaciones mínimas}
Primero se linealiza la población:
\begin{equation}
Y=\ln(P)
\end{equation}
donde $Y$ es la variable transformada y $P$ es la población.

Luego se ajusta el modelo lineal:
\begin{equation}
Y=a+bx
\end{equation}
donde $a$ es la ordenada al origen, $b$ es la pendiente y $x$ es el tiempo en años.

Finalmente se regresa a población:
\begin{equation}
P=e^{a+bx}=e^ae^{bx}
\end{equation}

En Excel:
\[
\texttt{=LINEST(rango\_lnP, rango\_x, TRUE, TRUE)}
\]
En Excel en español:
\[
\texttt{=ESTIMACION.LINEAL(rango\_lnP; rango\_x; VERDADERO; VERDADERO)}
\]

\excelnote{Agregue una columna con \texttt{=LN(P)}. Después use \texttt{LINEST} o \texttt{ESTIMACION.LINEAL} sobre esa nueva columna. La proyección final se obtiene con \texttt{=EXP(a+b*x)}.}

\section{Procedimiento}
\begin{enumerate}
\item Construir la columna $\ln(P)$.
\item Ajustar una recta entre $x$ y $\ln(P)$.
\item Recuperar la ecuación exponencial con la función \texttt{EXP}.
\item Proyectar los años futuros.
\end{enumerate}

\section{Ejemplo resuelto}
Aplicando \texttt{LINEST} sobre $\ln(P)$, o su equivalente \texttt{ESTIMACION.LINEAL}, se obtiene aproximadamente:
\[
a=10.1655,\qquad b=0.013714\ \text{año}^{-1}
\]

Por tanto:
\[
\ln(P)=10.1655+0.013714x
\]

y el modelo en población es:
\[
P=e^{10.1655+0.013714x}
\]

Entonces:
\[
P_{2030}=e^{10.1655+0.013714(60)}=\num{59182}
\]
\[
P_{2040}=e^{10.1655+0.013714(70)}=\num{67881}
\]
\[
P_{2050}=e^{10.1655+0.013714(80)}=\num{77860}
\]

\resultbox{El ajuste lineal sobre $\ln(P)$ produce 59,182 habitantes en 2030, 67,881 en 2040 y 77,860 en 2050.}

\section{Teoría breve}
Este método y \texttt{LOGEST} responden a la misma idea de fondo, pero aquí el procedimiento queda completamente transparente: se linealiza primero y se ajusta después.

\section{Observaciones de uso}
\begin{itemize}
\item Es didácticamente muy útil porque deja visible la transformación logarítmica.
\item Facilita explicar por qué un ajuste exponencial puede tratarse con herramientas lineales.
\item Debe recordarse que el ajuste ocurre sobre $\ln(P)$, no directamente sobre $P$.
\end{itemize}
\end{document}
